\subsection{群胚}

定义 4. —— 群胚是每个箭都可逆的范畴。

设 G 为群胚，a 为 G 的顶点。幺半群 \(G_{a,a}\) 是一个群，记作 \(G_{a}\) ，称为 G 在 a 处的迷向群。

设 \(a\) 与 \(b\) 为 G 的顶点，\(f \in \mathrm{G}_{a,b}\) 为连接 \(a\) 与 \(b\) 的箭。映射 \(g \mapsto fgf^{-1}\) 是从 \(\mathrm{G}_b\) 到 \(\mathrm{G}_a\) 的群同构，记作 \(\operatorname{Int}(f)\) 。若 \(f\) 与 \(f'\) 是 G 的可复合箭，则有 \(\operatorname{Int}(ff') = \operatorname{Int}(f) \circ \operatorname{Int}(f')\) 。同构 \(\operatorname{Int}(f)\) 的逆是同构 \(\operatorname{Int}(f^{-1})\) 。

群胚态射 \(\varphi\colon G\to G'\) 是范畴态射。此外，\(\varphi\) 是群胚同构当且仅当它是范畴同构。

设 \(\varphi\colon\mathrm{G}\to\mathrm{G}^{\prime}\) 为群胚态射。对 G 的每个顶点 \(a\) ，映射 \(f\mapsto\varphi(f)\) （从 \(\mathrm{G}_{a}\) 到 \((\mathrm{G}^{\prime})_{\varphi(a)}\) 的）是群同态，记作 \(\varphi_{a}\) 。特别地，有 \(\varphi(e_{a})=e_{\varphi(a)}\) 。

对 G 的每个箭 \(f\) ，\(\varphi(f^{-1})\) 是 \(\varphi(f)\) 的逆。若 \(f\) 是 G 中连接顶点 \(a\) 与顶点 \(b\) 的箭，则对每个元素 \(g \in \mathrm{G}_b\) ，有关系 \(\operatorname{Int}(\varphi(f))(\varphi(g)) = \varphi(\operatorname{Int}(f)(g))\) 。

